\chapter{Hamiltonian Truncation, DLCQ, And Benchmark Protocols}
\label{ch:constructive-hamiltonian-truncation-dlcq-benchmarks}

\ConventionsLorentzian

Hamiltonian numerical methods are regulator frameworks for QFT.  This
chapter formulates Hamiltonian truncation, light-front discretization, and
benchmark scripts as controlled approximations with explicitly declared
Hilbert spaces, cutoffs, counterterms, and convergence tests.

\section{Projected Hamiltonian Datum}

Let \(H\) be a self-adjoint Hamiltonian on a Hilbert space \(\Hilb\), bounded
below.  Choose finite-dimensional subspaces
\[
  \Hilb_\Lambda\subset\Hilb,\qquad P_\Lambda:\Hilb\to\Hilb_\Lambda,
\]
with \(P_\Lambda\to I\) strongly.  The projected Hamiltonian is
\[
  H_\Lambda=P_\Lambda H P_\Lambda
  \]
on \(\Hilb_\Lambda\).  If \(H\) is defined by a QFT regulator, \(H_\Lambda\)
must be supplemented by local counterterms:
\[
  H_\Lambda^{\mathrm{ren}}
  =
  P_\Lambda H P_\Lambda
  +
  \sum_a c_a(\Lambda)P_\Lambda\mathcal O_aP_\Lambda.
\]
The pair \((\Hilb_\Lambda,H_\Lambda^{\mathrm{ren}})\) is the numerical
regulator datum.

\section{Rayleigh--Ritz Bound}

Assume \(H\) is self-adjoint and has discrete eigenvalues
\[
  E_0\le E_1\le\cdots
\]
below an essential spectral threshold.  The \(n\)-th variational eigenvalue
in a finite subspace is
\[
  E_n(\Lambda)
  =
  \min_{\dim W=n+1,\ W\subset\Hilb_\Lambda}
  \max_{\psi\in W,\ \|\psi\|=1}\langle\psi,H\psi\rangle.
\]
Since the minimization is over fewer subspaces than in the full Hilbert
space,
\[
  E_n(\Lambda)\ge E_n.
\]
If \(\Hilb_\Lambda\) is nested and dense in the form domain, these variational
eigenvalues decrease to \(E_n\).  In QFT, counterterms and continuum limits
can destroy this monotonicity unless the comparison is made at fixed
regulated Hamiltonian.

\section{Truncated Conformal Space}

For a two-dimensional QFT built as a relevant deformation of a CFT, choose the
CFT Hilbert space on a circle and keep states satisfying
\[
  L_0+\bar L_0-\frac{c}{12}\le E_{\mathrm{cut}}\,\frac{L}{2\pi}.
\]
The Hamiltonian matrix elements are
\[
  \langle i|H|j\rangle
  =
  \frac{2\pi}{L}
  \left(\Delta_i-\frac{c}{12}\right)\delta_{ij}
  +
  \sum_a\lambda_a
  \left(\frac{L}{2\pi}\right)^{1-\Delta_a}
  \langle i|\mathcal O_a(1)|j\rangle.
\]
High-energy states generate cutoff dependence.  The local Hamiltonian
counterterms are determined by matching low-energy observables as
\(E_{\mathrm{cut}}\) is varied.

\section{DLCQ Datum}

Light-front coordinates are
\[
  x^\pm=\frac1{\sqrt2}(x^0\pm x^1).
\]
DLCQ compactifies
\[
  x^-\sim x^-+2\pi R
\]
so that
\[
  p^+=\frac{n}{R},\qquad n\in\mathbb Z_{\ge0}.
\]
The total longitudinal momentum is
\[
  P^+=\frac{K}{R},
\]
where \(K\) is the harmonic resolution.  The finite DLCQ Hilbert space at
fixed \(K\) consists of partitions of \(K\) among partons, with transverse and
internal cutoffs added as needed.  Zero modes with \(n=0\) are constrained
variables or additional degrees of freedom depending on the theory; their
treatment is part of the regulator.

\section{Light-Front Hamiltonian}

The invariant mass operator is
\[
  M^2=2P^+P^- -P_\perp^2.
\]
At fixed \(P^+\) and \(P_\perp\), diagonalizing \(P^-\) is equivalent to
diagonalizing \(M^2\).  Gauge theories require solving constraint equations
for nondynamical fields and preserving residual gauge constraints.  A
light-front truncation without the constraint algebra and zero-mode sector is
not a definition of the gauge theory.

\section{Benchmark Protocol}

A benchmark for any truncation method should specify:
\[
\begin{gathered}
  \text{regulated Hilbert space},\qquad
  \text{operator basis},\qquad
  \text{counterterm coordinates},\\
  \text{observables used for matching},\qquad
  \text{extrapolation model},\qquad
  \text{independent checks}.
\end{gathered}
\]
Independent checks include Ward identities, known finite-volume limits,
reflection positivity where applicable, spectral sum rules, and comparison
with Euclidean lattice or exact integrable data in overlapping regimes.

\begin{openproblem}[Certified QFT truncation]
\label{op:certified-qft-hamiltonian-truncation}
Develop Hamiltonian truncation and DLCQ schemes with computable error bounds
for interacting QFT observables, including counterterm determination,
constraint preservation, zero-mode treatment, continuum extrapolation, and
public benchmark scripts.
\end{openproblem}
